\documentclass[10pt]{amsart}
\usepackage[a4paper,margin=23mm]{geometry}
\usepackage{amsmath,amssymb,amsthm,mathtools,hyperref,graphicx,xspace,letltxmacro,xpatch}
\usepackage[english]{babel}
\newtheorem{theo}{Theorem}[section]
\newtheorem{lemm}[theo]{Lemma}\newtheorem{prop}[theo]{Proposition}\newtheorem{coro}[theo]{Corollary}\newtheorem{conj}[theo]{Conjecture}
\theoremstyle{definition}\newtheorem{defi}[theo]{Definition}\newtheorem{rema}[theo]{Remark}\newtheorem{exam}[theo]{Example}\newtheorem{exem}[theo]{Example}\newtheorem{nota}[theo]{Notation}
\newcommand{\ie}{i.e.}\newcommand{\eg}{e.g.}\newcommand{\resp}{resp.}\newcommand{\skpt}{\par\smallskip\noindent}
\emergencystretch=3em\pagestyle{plain}
%~Mouliné par MaN_auto v.0.23.3 2021-05-26 11:22:15

\usepackage[refpage,noprefix,notocbasic]{nomencl}
\usepackage{bm}
%\usepackage{stmaryrd}
\usepackage[lofdepth,lotdepth]{subfig}
%\usepackage{paralist}





%
\AtBeginDocument{
\newtheorem{assumption}{Assumption}
}
\newcommand*\theassumption{\arabic{assumption}}



\providecommand{\R}{}
\providecommand{\Z}{}
\providecommand{\N}{}
\providecommand{\C}{}
\providecommand{\Q}{}
\providecommand{\G}{}
\providecommand{\Lt}{}

\renewcommand{\R}{\mathbb{R}}
\renewcommand{\N}{\mathbb{N}}
\newcommand{\bbZ}{\mathbb{Z}}
\newcommand{\bbP}{\mathbb{P}}
\newcommand{\bbT}{\mathbb{T}}

\newcommand{\frkS}{\mathfrak{S}}


\newcommand{\bfe}{\mathbf{e}}
\newcommand{\be}{\mathbf{e}}
\newcommand{\bfq}{\mathbf{q}}


\newcommand{\cL}{\mathcal{L}}
\newcommand{\cM}{\mathcal{M}}
\newcommand{\cT}{\mathcal{T}}
\newcommand{\cZ}{\mathcal{Z}}
\newcommand{\cC}{\mathcal{C}}
\newcommand{\cO}{\mathcal{O}}



\newcommand{\rD}{\mathrm{D}}
\newcommand{\rQ}{\mathrm{Q}}
\newcommand{\rR}{\mathrm{R}}
\newcommand{\rT}{\mathrm{T}}
\newcommand{\rX}{\mathrm{X}}
\newcommand{\rXi}{\mathrm{X'}}


\newcommand{\bM}{\mathbf{M}}
\newcommand{\bT}{\mathbf{T}}
\newcommand{\bP}{\mathbf{P}}
\newcommand{\bE}{\mathbf{E}}


\newcommand{\rv}{\mathrm{v}}
\newcommand{\rd}{\mathrm{d}}
\newcommand{\rr}{\mathrm{r}}
\newcommand{\rs}{\mathrm{s}}
\newcommand{\rz}{\mathrm{z}}
\newcommand{\ru}{\mathrm{u}}
\newcommand{\rt}{\mathrm{t}}
\newcommand{\rp}{\mathrm{p}}


\DeclareMathOperator{\sym}{sym}
\DeclareMathOperator{\dist}{dist}
\DeclareMathOperator{\Dist}{Dist}
\DeclareMathOperator{\branch}{branch }
\DeclareMathOperator{\cht}{ctype}
\DeclareMathOperator{\Deg}{deg}
\DeclareMathOperator{\fin}{fin}
\DeclareMathOperator{\height}{height}


%\providecommand{deg}{}


\newcommand{\lip}[1]{\ensuremath{|\!|#1|\!|_{\text{Lip}}}}
\newcommand{\ninf}[1]{\ensuremath{|\!|#1|\!|_{\infty}}}

\newcommand{\E}[1]{\mathbf{E}\left[#1\right]}
\newcommand{\p}[1]{\mathbf{P}\left\{#1\right\}}
\newcommand{\set}[1]{\left\{ #1 \right\}}

\newcommand{\probC}[2]{\mathbf{P}\set{#1 \; \left| \; #2 \right.}}

\newcommand{\limC}{\ensuremath{C}}
\newcommand{\limZ}{\ensuremath{Z}}
\newcommand{\sbt}{{\ensuremath{\scriptscriptstyle\bullet}}}

\newcommand{\Permvecs}{P}
\newcommand{\unif}{\ensuremath{X}}
\newcommand{\dfun}{\ensuremath{\delta}}



\newcommand{\angles}[1]{\langle #1 \rangle}

\providecommand{\obar}{}
\providecommand{\eps}{}
\renewcommand{\eps}{\epsilon}
\providecommand{\ora}[1]{}


\newcommand{\llbracket}{\mathopen{[\mkern -2.7mu[}}
\newcommand{\rrbracket}{\mathclose{]\mkern -2.7mu]}}
%\newcommand{\bbr}[1]{\llbracket{ #1 \rrbracket}



\newcommand{\eqdist}{\ensuremath{\stackrel{\mathrm{d}}{=}}}
\newcommand{\convdist}{\ensuremath{\stackrel{\mathrm{d}}{\rightarrow}}}
\newcommand{\convpr}{\ensuremath{\stackrel{\mathrm{(p)}}{\longrightarrow}}}
\newcommand{\aseq}{\ensuremath{\stackrel{\mathrm{a.s.}}{=}}}


\renewcommand*{\pagedeclaration}[1]{\unskip\dotfill\hyperpage{#1}}
\makenomenclature
\setlength{\nomitemsep}{-\parsep}



%\DeclareRobustCommand{\SkipTocEntry}[5]{}




%%%%%%%%%%%%%%%%%%%%%%%%%%%%%%%%%%%%%%%%%%%%%%%%%%%%%%%%%%%

\graphicspath{{./figures/}}

\newcommand*{\mk}{\mkern -1mu}
\newcommand*{\Mk}{\mkern -2mu}
\newcommand*{\mK}{\mkern 1mu}
\newcommand*{\MK}{\mkern 2mu}

%\hypersetup{urlcolor=purple, linkcolor=blue, citecolor=red}

\newcommand*{\romanenumi}{\renewcommand*{\theenumi}{\roman{enumi}}}
\newcommand*{\Romanenumi}{\renewcommand*{\theenumi}{\Roman{enumi}}}
\newcommand*{\alphenumi}{\renewcommand*{\theenumi}{\alph{enumi}}}
\newcommand*{\Alphenumi}{\renewcommand*{\theenumi}{\Alph{enumi}}}
\let\oldtilde\tilde
\renewcommand*{\tilde}[1]{\mathchoice{\widetilde{#1}}{\widetilde{#1}}{\oldtilde{#1}}{\oldtilde{#1}}}
\let\oldhat\hat
\renewcommand*{\hat}[1]{\mathchoice{\widehat{#1}}{\widehat{#1}}{\oldhat{#1}}{\oldhat{#1}}}
\let\oldforall\forall
\renewcommand*{\forall}{\mathrel{\oldforall}}
%%%%%%%%%%%%%%%%%%%%%%%%%%%%%%%%%%%%%%%%%%%%%%%%%%%%%%%%%%%%


\NewDocumentEnvironment{enonce*}{ O{} m O{} }{\par\smallskip\noindent\textbf{#2\if\relax\detokenize{#3}\relax\else\ (#3)\fi.}\ }{\par\smallskip}
\NewDocumentEnvironment{enonce}{ O{} m O{} }{\par\smallskip\refstepcounter{theo}\noindent\textbf{#2~\thetheo\if\relax\detokenize{#3}\relax\else\ (#3)\fi.}\ }{\par\smallskip}
\newenvironment{proofc}[1][\proofname]{\begin{proof}[{#1}]}{\end{proof}}
